\documentclass{article}
\usepackage[utf8]{inputenc}
\usepackage{amsmath}
\usepackage{amsfonts}
\usepackage{amssymb}
\usepackage{graphicx}
\usepackage{hyperref}
\usepackage{geometry}
\usepackage{parskip}
\geometry{a4paper, margin=1in}

\title{TensorStore: An io\_uring-Native Evolution of the ServerlessLLM Storage Engine\\
Project Plan}
\author{Botir Khaltaev\\
CS3821 - Final Year Project\\
Supervised By: Daniel O'Keeffe\\
Department of Computer Science\\
Royal Holloway, University of London}
\date{}

\begin{document}

\maketitle

\section{Abstract}

Large Language Model (LLM) deployment encounters a critical performance bottleneck: loading times exceeding several minutes for multi-billion parameter models. While ServerlessLLM (Fu et al., 2024) demonstrated 5-10x loading improvements through multi-threaded pipelines, this approach incurs substantial CPU overhead from context switching and synchronization primitives.

This project investigates whether io\_uring, Linux's high-performance asynchronous I/O interface, can exceed multi-threaded performance by eliminating these architectural inefficiencies. TensorStore proposes an io\_uring-native event loop architecture designed to maximize I/O throughput while minimizing CPU utilization.

The research pursues two primary objectives: first, empirically validating io\_uring's performance advantages over traditional tokio::fs approaches for tensor loading; second, implementing a custom TensorStore format specifically optimized for io\_uring's vectored operations and batching capabilities.

Success criteria include achieving >20\% performance improvement over baseline implementations, supported by comprehensive empirical analysis identifying the conditions under which these benefits materialize. Within the 300-hour timeframe, the project prioritizes rapid hypothesis validation, with conditional advancement to production features contingent upon demonstrated performance gains.

\section{Timeline}


\subsection{Phase 1: Research (Weeks 1-2, 40 hours)}

\begin{itemize}
\item Analyzed ServerlessLLM architecture and performance bottlenecks
\item Researched io\_uring ecosystem and tokio-uring bindings
\item Designed TensorStore format specification for vectored I/O
\item Created comprehensive research documentation
\end{itemize}

\subsection{Phase 2: Core Validation (Week 3, 20 hours)}

\begin{itemize}
\item \textbf{Setup (2h):} Rust project with tokio-uring, safetensors, criterion
\item \textbf{Baseline (6h):} Safetensors loading with tokio::fs
\item \textbf{Test (6h):} Safetensors loading with tokio-uring
\item \textbf{TensorStore (6h):} Custom format exploiting vectored I/O
\end{itemize}

\subsection{Phase 3: Analysis \& Documentation (Week 4, 30 hours)}

\begin{itemize}
\item \textbf{Performance Analysis (20h):} CPU/memory profiling, bottleneck identification
\item \textbf{Documentation (10h):} Comprehensive findings report
\end{itemize}

\subsection{Phase 4: Advanced Features (Weeks 5-15, 210 hours)}

\textbf{Core Features:}
\begin{itemize}
\item \textbf{Vectored I/O (50h):} IORING\_OP\_READV, batch operations
\item \textbf{Format Enhancement (40h):} Advanced TensorStore features
\end{itemize}

\textbf{Production Features (if >20\% improvement):}
\begin{itemize}
\item \textbf{Framework Integration (60h):} PyTorch, Python bindings
\item \textbf{NUMA \& Multi-GPU (30h):} Hardware optimization
\item \textbf{Production Polish (30h):} Error handling, documentation
\end{itemize}

\section{Risks and Mitigations}

Systems-level research involving nascent asynchronous I/O technologies necessitates comprehensive risk management across multiple dimensions.

\subsection{Hardware Failure}

Local hardware failures pose risks to project continuity and data integrity. Mitigation employs continuous GitHub version control with frequent commits, coupled with cloud-based documentation infrastructure (Overleaf), ensuring complete work preservation independent of local system failures.

\subsection{Time Estimation Accuracy}

The fixed 300-hour constraint demands precise task estimation. Risk mitigation strategies include granular task decomposition, a critical Week 3 go/no-go checkpoint, and optional implementation phases providing buffer capacity for unforeseen complexities.

\subsection{Hypothesis Invalidation}

Negative results—wherein io\_uring fails to demonstrate anticipated performance advantages—constitute valuable scientific contributions. This risk is addressed through:
\begin{itemize}
\item Clearly defined success thresholds (>20\% improvement benchmark)
\item Rigorous root-cause analysis of performance characteristics
\item Systematic documentation of negative findings as legitimate research outcomes
\item TensorStore format utility independent of performance hypothesis validation
\end{itemize}

\subsection{Ecosystem Maturity}

The tokio-uring crate exhibits limited recent development activity, introducing potential stability and compatibility vulnerabilities. Mitigation strategies encompass:
\begin{itemize}
\item Comprehensive early-phase integration testing
\item Identification and preparation of alternative io\_uring binding libraries
\item Systematic documentation of encountered ecosystem limitations
\item Strategic upstream contributions to address critical deficiencies
\end{itemize}

\subsection{Platform Constraints}

io\_uring's Linux-specific implementation fundamentally constrains solution portability. This limitation is addressed through:
\begin{itemize}
\item Deliberate focus on Linux deployment environments where io\_uring excels
\item Comprehensive architectural documentation facilitating cross-platform adaptation
\item Generalizable methodology applicable to alternative asynchronous I/O interfaces
\item Explicit acknowledgment of platform dependencies within research findings
\end{itemize}

\subsection{Measurement Validity}

Accurate performance characterization in systems research demands rigorous experimental methodology. The measurement framework incorporates:
\begin{itemize}
\item Industry-standard Criterion benchmarking infrastructure
\item Comprehensive test matrices spanning diverse file sizes and access patterns
\item Multi-dimensional profiling encompassing CPU utilization, memory allocation, and I/O characteristics
\item Detailed hardware configuration documentation ensuring reproducibility
\item Statistical validation through extensive measurement repetition
\end{itemize}

\subsection{Implementation Complexity}

io\_uring's sophisticated API presents substantial implementation challenges, with poorly designed code potentially underperforming baseline approaches. Risk mitigation employs:
\begin{itemize}
\item Systematic complexity escalation from foundational to advanced features
\item Strict adherence to established tokio-uring design patterns and best practices
\item Comprehensive error handling and edge case management
\item Continuous performance regression detection through incremental validation
\end{itemize}

\section{Success Criteria}

Rigorous success metrics establish project value independent of hypothesis validation outcomes, ensuring meaningful research contributions under all experimental conditions.

\subsection{Primary Success Criteria}

\begin{itemize}
\item \textbf{Performance Validation:} Measurable comparison between tokio-uring and tokio::fs
\item \textbf{Technical Feasibility:} Working io\_uring tensor loading implementation
\item \textbf{Empirical Analysis:} Documentation of performance factors and conditions
\item \textbf{Research Contribution:} Actionable insights for ML workload optimization
\end{itemize}

\subsection{Secondary Success Criteria}

\begin{itemize}
\item \textbf{TensorStore Implementation:} Custom format demonstrating I/O optimization
\item \textbf{Framework Integration:} PyTorch integration (if hypothesis succeeds)
\item \textbf{Production Features:} Vectored I/O and NUMA support (if benefits proven)
\item \textbf{Documentation:} Technical documentation for researcher adoption
\end{itemize}

\subsection{Success Metrics}

\begin{itemize}
\item \textbf{Performance Target:} >20\% improvement over tokio::fs baseline
\item \textbf{Code Quality:} Robust implementation with error handling
\item \textbf{Analysis Depth:} Comprehensive profiling data
\item \textbf{Decision Clarity:} Clear adoption recommendation
\end{itemize}

This research framework guarantees substantial contributions to the field regardless of empirical outcomes, thereby optimizing the value derived from the allocated 300-hour investigation period.

\begin{thebibliography}{9}

\bibitem{serverlessllm2024}
Fu, Y., Xue, L., Huang, Y., Brabete, A., Ustiugov, D., Patel, Y. and Mai, L. (2024) 'ServerlessLLM: Low-Latency Serverless Inference for Large Language Models', \textit{18th USENIX Symposium on Operating Systems Design and Implementation (OSDI 24)}. Santa Clara, CA: USENIX Association, pp. 135-153. Available at: https://www.usenix.org/conference/osdi24/presentation/fu

\bibitem{tokio_blog}
Tokio Team (2021) \textit{Announcing io\_uring support for Tokio}. Available at: https://tokio.rs/blog/2021-07-tokio-uring

\bibitem{tokio_uring}
Tokio Contributors (2025) \textit{tokio-uring: io\_uring support for Tokio}. GitHub. Available at: https://github.com/tokio-rs/tokio-uring

\bibitem{safetensors}
HuggingFace (2024) \textit{Safetensors: Simple, safe way to store and distribute tensors}. GitHub. Available at: https://github.com/huggingface/safetensors

\bibitem{iouring_intro}
Axboe, J. (2019) \textit{Efficient IO with io\_uring}. Available at: https://kernel.dk/io\_uring.pdf

\bibitem{liburing}
Axboe, J. (2025) \textit{liburing: io\_uring library}. GitHub. Available at: https://github.com/axboe/liburing

\bibitem{serverlessllm_repo}
ServerlessLLM Contributors (2024) \textit{ServerlessLLM: Locality-enhanced Serverless Inference for Large Language Models}. GitHub. Available at: https://github.com/ServerlessLLM/ServerlessLLM

\end{thebibliography}

\end{document}
